% Options for packages loaded elsewhere
\PassOptionsToPackage{unicode}{hyperref}
\PassOptionsToPackage{hyphens}{url}
\documentclass[
]{book}
\usepackage{xcolor}
\usepackage{amsmath,amssymb}
\setcounter{secnumdepth}{5}
\usepackage{iftex}
\ifPDFTeX
  \usepackage[T1]{fontenc}
  \usepackage[utf8]{inputenc}
  \usepackage{textcomp} % provide euro and other symbols
\else % if luatex or xetex
  \usepackage{unicode-math} % this also loads fontspec
  \defaultfontfeatures{Scale=MatchLowercase}
  \defaultfontfeatures[\rmfamily]{Ligatures=TeX,Scale=1}
\fi
\usepackage{lmodern}
\ifPDFTeX\else
  % xetex/luatex font selection
\fi
% Use upquote if available, for straight quotes in verbatim environments
\IfFileExists{upquote.sty}{\usepackage{upquote}}{}
\IfFileExists{microtype.sty}{% use microtype if available
  \usepackage[]{microtype}
  \UseMicrotypeSet[protrusion]{basicmath} % disable protrusion for tt fonts
}{}
\makeatletter
\@ifundefined{KOMAClassName}{% if non-KOMA class
  \IfFileExists{parskip.sty}{%
    \usepackage{parskip}
  }{% else
    \setlength{\parindent}{0pt}
    \setlength{\parskip}{6pt plus 2pt minus 1pt}}
}{% if KOMA class
  \KOMAoptions{parskip=half}}
\makeatother
\usepackage{longtable,booktabs,array}
\usepackage{calc} % for calculating minipage widths
% Correct order of tables after \paragraph or \subparagraph
\usepackage{etoolbox}
\makeatletter
\patchcmd\longtable{\par}{\if@noskipsec\mbox{}\fi\par}{}{}
\makeatother
% Allow footnotes in longtable head/foot
\IfFileExists{footnotehyper.sty}{\usepackage{footnotehyper}}{\usepackage{footnote}}
\makesavenoteenv{longtable}
\usepackage{graphicx}
\makeatletter
\newsavebox\pandoc@box
\newcommand*\pandocbounded[1]{% scales image to fit in text height/width
  \sbox\pandoc@box{#1}%
  \Gscale@div\@tempa{\textheight}{\dimexpr\ht\pandoc@box+\dp\pandoc@box\relax}%
  \Gscale@div\@tempb{\linewidth}{\wd\pandoc@box}%
  \ifdim\@tempb\p@<\@tempa\p@\let\@tempa\@tempb\fi% select the smaller of both
  \ifdim\@tempa\p@<\p@\scalebox{\@tempa}{\usebox\pandoc@box}%
  \else\usebox{\pandoc@box}%
  \fi%
}
% Set default figure placement to htbp
\def\fps@figure{htbp}
\makeatother
\setlength{\emergencystretch}{3em} % prevent overfull lines
\providecommand{\tightlist}{%
  \setlength{\itemsep}{0pt}\setlength{\parskip}{0pt}}
\usepackage[]{natbib}
\bibliographystyle{plainnat}
\usepackage{booktabs}
\usepackage{bookmark}
\IfFileExists{xurl.sty}{\usepackage{xurl}}{} % add URL line breaks if available
\urlstyle{same}
\hypersetup{
  pdftitle={Ciência de dados e Big Data},
  pdfauthor={Professor Pedro Luiz Pzzigatti Correa},
  hidelinks,
  pdfcreator={LaTeX via pandoc}}

\title{Ciência de dados e Big Data}
\author{Professor Pedro Luiz Pzzigatti Correa}
\date{2026-09-24}

\begin{document}
\maketitle

{
\setcounter{tocdepth}{1}
\tableofcontents
}
\chapter*{Sobre estas anotações}\label{sobre-estas-anotauxe7uxf5es}
\addcontentsline{toc}{chapter}{Sobre estas anotações}

---------------------------------------------------------------------------------------------------------------------------------------

Estas anotações são apenas lembretes das aulas expostas em sala, durante a disciplina de Data Science e Big Data.

\section*{Calendário das aulas}\label{calenduxe1rio-das-aulas}
\addcontentsline{toc}{section}{Calendário das aulas}

---------------------------------------------------------------------------------------------------------------------------------------

\paragraph*{SETEMBRO DE 2026}\label{setembro-de-2026}
\addcontentsline{toc}{paragraph}{SETEMBRO DE 2026}

\begin{longtable}[]{@{}
  >{\raggedright\arraybackslash}p{(\linewidth - 6\tabcolsep) * \real{0.0813}}
  >{\raggedright\arraybackslash}p{(\linewidth - 6\tabcolsep) * \real{0.1057}}
  >{\raggedright\arraybackslash}p{(\linewidth - 6\tabcolsep) * \real{0.0488}}
  >{\raggedright\arraybackslash}p{(\linewidth - 6\tabcolsep) * \real{0.7642}}@{}}
\toprule\noalign{}
\begin{minipage}[b]{\linewidth}\raggedright
Data
\end{minipage} & \begin{minipage}[b]{\linewidth}\raggedright
Dia da Semana
\end{minipage} & \begin{minipage}[b]{\linewidth}\raggedright
Aulas
\end{minipage} & \begin{minipage}[b]{\linewidth}\raggedright
Conteúdo
\end{minipage} \\
\midrule\noalign{}
\endhead
\bottomrule\noalign{}
\endlastfoot
24/09/2026 & Quinta-Feira & Aula 1 & Apresentação da Disciplina. Introdução a Modelos de Gestão de Dados. \textbf{Lista de Exercícios 1} \\
\end{longtable}

\paragraph*{OUTUBRO DE 2026}\label{outubro-de-2026}
\addcontentsline{toc}{paragraph}{OUTUBRO DE 2026}

\begin{longtable}[]{@{}
  >{\raggedright\arraybackslash}p{(\linewidth - 6\tabcolsep) * \real{0.0694}}
  >{\raggedright\arraybackslash}p{(\linewidth - 6\tabcolsep) * \real{0.0903}}
  >{\raggedright\arraybackslash}p{(\linewidth - 6\tabcolsep) * \real{0.0417}}
  >{\raggedright\arraybackslash}p{(\linewidth - 6\tabcolsep) * \real{0.7986}}@{}}
\toprule\noalign{}
\begin{minipage}[b]{\linewidth}\raggedright
Data
\end{minipage} & \begin{minipage}[b]{\linewidth}\raggedright
Dia da Semana
\end{minipage} & \begin{minipage}[b]{\linewidth}\raggedright
Aulas
\end{minipage} & \begin{minipage}[b]{\linewidth}\raggedright
Conteúdo
\end{minipage} \\
\midrule\noalign{}
\endhead
\bottomrule\noalign{}
\endlastfoot
01/10/2026 & Quinta-Feira & Aula 2 & Definição dos Grupos e Temas dos Projetos. \textbf{Atividade somente dos grupos -- definição dos temas} \\
08/10/2026 & Quinta-Feira & Aula 3 & Análise Exploratória dos Dados. \textbf{Lista de Exercícios 2} \\
15/10/2026 & Quinta-Feira & Aula 4 & Preparação dos Dados, Seleção e definição de amostras. Análise de Desempenho dos Modelos. \textbf{Lista de Exercícios 3} \\
22/10/2026 & Quinta-Feira & Aula 5 & Análise e Visualização dos Dados. \textbf{Lista de Exercícios 4} \\
29/10/2026 & Quinta-Feira & Aula 6 & Desenvolvimento do Projeto da Disciplina \\
\end{longtable}

\paragraph*{NOVEMBRO DE 2026}\label{novembro-de-2026}
\addcontentsline{toc}{paragraph}{NOVEMBRO DE 2026}

\begin{longtable}[]{@{}
  >{\raggedright\arraybackslash}p{(\linewidth - 6\tabcolsep) * \real{0.1587}}
  >{\raggedright\arraybackslash}p{(\linewidth - 6\tabcolsep) * \real{0.2063}}
  >{\raggedright\arraybackslash}p{(\linewidth - 6\tabcolsep) * \real{0.1111}}
  >{\raggedright\arraybackslash}p{(\linewidth - 6\tabcolsep) * \real{0.5238}}@{}}
\toprule\noalign{}
\begin{minipage}[b]{\linewidth}\raggedright
Data
\end{minipage} & \begin{minipage}[b]{\linewidth}\raggedright
Dia da Semana
\end{minipage} & \begin{minipage}[b]{\linewidth}\raggedright
Aulas
\end{minipage} & \begin{minipage}[b]{\linewidth}\raggedright
Conteúdo
\end{minipage} \\
\midrule\noalign{}
\endhead
\bottomrule\noalign{}
\endlastfoot
05/11/2026 & Quinta-Feira & Aula 7 & Seminário -- Projeto da Disciplina \\
12/11/2026 & Quinta-Feira & Aula 8 & Seminário -- Projeto da Disciplina \\
19/11/2026 & Quinta-Feira & Aula 9 & Seminário -- Projeto da Disciplina \\
26/11/2026 & Quinta-Feira & Aula 10 & Seminário -- Projeto da Disciplina \\
\end{longtable}

\paragraph*{DEZEMBRO DE 2026}\label{dezembro-de-2026}
\addcontentsline{toc}{paragraph}{DEZEMBRO DE 2026}

\begin{longtable}[]{@{}
  >{\raggedright\arraybackslash}p{(\linewidth - 6\tabcolsep) * \real{0.1493}}
  >{\raggedright\arraybackslash}p{(\linewidth - 6\tabcolsep) * \real{0.1940}}
  >{\raggedright\arraybackslash}p{(\linewidth - 6\tabcolsep) * \real{0.1045}}
  >{\raggedright\arraybackslash}p{(\linewidth - 6\tabcolsep) * \real{0.5522}}@{}}
\toprule\noalign{}
\begin{minipage}[b]{\linewidth}\raggedright
Data
\end{minipage} & \begin{minipage}[b]{\linewidth}\raggedright
Dia da Semana
\end{minipage} & \begin{minipage}[b]{\linewidth}\raggedright
Aulas
\end{minipage} & \begin{minipage}[b]{\linewidth}\raggedright
Conteúdo
\end{minipage} \\
\midrule\noalign{}
\endhead
\bottomrule\noalign{}
\endlastfoot
03/12/2026 & Quinta-Feira & Aula 11 & \textbf{Data Final de entrega dos Artigos} \\
\end{longtable}

\chapter{- Introdução a Modelos de Gestão de Dados}\label{introduuxe7uxe3o-a-modelos-de-gestuxe3o-de-dados}

\chapter{- Definição de Grupos e Temas de Projetos}\label{definiuxe7uxe3o-de-grupos-e-temas-de-projetos}

\chapter{- Análise Exploratória de dados}\label{anuxe1lise-exploratuxf3ria-de-dados}

\chapter{- Preparação de Dados, Seleção e Definição de Amostras, Análsie de Desempenho de Modelos}\label{preparauxe7uxe3o-de-dados-seleuxe7uxe3o-e-definiuxe7uxe3o-de-amostras-anuxe1lsie-de-desempenho-de-modelos}

\bibliography{book.bib}

\end{document}
